% =====================================================================
% DRAFT (10 Oct 2026) -- Section: computational cost
% Every number carries a "% fonte:" comment. Missing numbers are \todo{}.
% Main sources: aau/runs/indomain_recog/summary.md ("Tempi", job 1060419,
% one L40S node, CPU parts on ONE thread) and aau/runs/evidence/e10/E10.md
% (DiffusionNet operators on GPU, V100, jobs 1061877/1061896/1061897).
% =====================================================================

\section{Computational Cost}
\label{sec:cost}

Under FR, pairwise non-rigid registration remains more accurate than our
current encoder on HIFI3D (Table~\ref{tab:gt_reeval}). The case for a learned
metric therefore rests partly on cost, and we measure it rather than
estimate it wherever possible. The two families of
Sect.~\ref{sec:protocol_baselines} scale differently. A pairwise pipeline
pays a registration for every comparison, $O(P\,C_{\mathrm{reg}})$ for $P$
pairs. An enrolment method pays once per mesh and then compares fixed-size
descriptors, $O(N\,C_{\mathrm{enr}}+P\,C_{\mathrm{cmp}})$ for $N$ meshes,
with $C_{\mathrm{cmp}}\ll C_{\mathrm{reg}}$. For identification of one query
against a gallery of $N$ enrolled meshes (1:$N$ search), the first costs $N$
registrations and the second one enrolment plus $N$ descriptor comparisons.

\paragraph{Setting.}
Unless stated otherwise, all timings were measured in a single job on one
node with an NVIDIA L40S GPU and an AMD EPYC 9454 CPU, with every CPU stage
restricted to one thread, as the pipelines run in practice, on 60 meshes
(5 BFM and 5 ICT subjects in six topologies, 3{,}275 to 60{,}435 vertices,
median 12{,}955) and 60 cross-topology pairs. We report medians and
interquartile ranges; meshes were read from RAM.
% fonte: aau/runs/indomain_recog/summary.md, sez. "Tempi" (protocollo) e intestazione delle tabelle (host a768-l40s-06, EPYC 9454, L40S, job 1060419, OMP_NUM_THREADS=1, 60 mesh 3275-60435 vertici, mediana 12955)
The timed encoder is the BFM+ICT checkpoint with the same backbone as the
models of this paper (256-dimensional embedding)
\todo{rimisurare iscrizione e ricerca con il modello finale (testa fattorizzata), stesso nodo, end-to-end}.
% fonte: aau/runs/indomain_recog/summary.md, sez. "Tempi" ("congiunto ... L2 su 256-d")

\paragraph{Enrolment.}
Enrolling a mesh with our encoder requires the DiffusionNet
operators~\cite{diffusionet} (Laplacian, mass, $k=128$ eigenpairs, gradients)
and one forward pass. On one CPU thread the operators dominate: 2.03\,s
(IQR 1.23--4.33\,s), against 55\,ms (IQR 39--94\,ms) for loading and the
forward pass on the GPU.
% fonte: aau/runs/indomain_recog/summary.md, "iscrizione (per mesh)" (operatori CPU 1 thread 2.03 s [1.23-4.33]; embedding GPU 54.96 ms [39.45-94.24])
We therefore moved the operator computation to the GPU: the eigenpairs are
obtained by shift-invert with a sparse Cholesky factorization and a block
Krylov iteration, which is exact to the CPU solution (embeddings within
$1.4\times10^{-6}$, the same order as the run-to-run noise of the CPU solver
itself, $1.2\times10^{-6}$).
% fonte: aau/runs/evidence/e10/E10.md, "Cosa fa la GPU" (via bk) e "Correttezza" (dz max 1.31e-6 / 1.40e-6 per k 128 / 64; cpu_v0 1.14e-6 / 1.22e-6)
On a V100 this takes 97\,ms per mesh for a 9.4k-vertex ICT mesh at $k=128$
(63\,ms at 3.3k vertices, 167\,ms at 24.1k, 309\,ms at 60.4k), against about
0.64\,s for one CPU process.
% fonte: aau/runs/evidence/e10/E10.md, "Per classe" (latenza B=1, k 128: 63 / 97 / 167 / 309 ms) e "Latenza" (~0.64 s CPU con grad_vec)
Operators and forward pass together therefore cost on the order of 0.15\,s per
mesh on a GPU; this figure combines two measurements on different GPUs (V100
for the operators, L40S for the forward pass) and different mesh samples, and
is not an end-to-end measurement
\todo{misura end-to-end operatori GPU + forward sullo stesso nodo L40S; E10 non ha misurato la L40S (fp64 a 1/64)}.
% fonte: aau/runs/evidence/e10/E10.md, "GPU non misurate: L40S" e Verdetto 1 ("Non misurato su L40S: la via e' fp64")
The strongest enrolment baseline, non-rigid registration to a domain template
(similarity ICP, NICP and a final Procrustes step), costs 1.15\,s per mesh
(IQR 1.13--1.31\,s) on one CPU thread.
% fonte: aau/runs/indomain_recog/summary.md, "iscrizione (per mesh)" (NICP su template 1.15 s [1.13-1.31])
On HIFI3D, with a 9{,}518-vertex template, the same enrolment took a median
of 2.79\,s per mesh in a separate run.
% fonte: aau/runs/competitors_hifi3d/summary.md, Controlli ("NICP su template ... 9518 vertici, 4096 tenuti ... mediana 2.79 s per mesh")
\todo{verificare hardware e numero di thread di quel run prima di citarlo}
ArcFace on three rendered views costs 0.57\,s per mesh on one CPU thread, of
which 0.23\,s is framing and rendering.
% fonte: aau/runs/indomain_recog/summary.md, "iscrizione (per mesh)" (ArcFace migliore 570.08 ms; inquadratura + 3 render 228.71 ms)
We did not port the template registration to the GPU, so the comparison of
enrolment times is between a GPU implementation of our pipeline and
single-thread CPU implementations of the baselines; on one CPU thread our
enrolment (2.08\,s) is slower than NICP on template.
% fonte: aau/runs/indomain_recog/summary.md, "iscrizione (per mesh)" (congiunto totale operatori CPU + embedding GPU 2.08 s); paper/PLAN_MASSIVE.md, sez. 16.6 ("E10 passa con riserve sui confronti di velocita'")

\paragraph{Search.}
Once enrolled, a comparison in our embedding space is a Euclidean distance
between two 256-dimensional vectors (2.4\,$\mu$s), and a 1:10{,}000 search
including sorting takes 1.96\,ms on one CPU thread and 68\,$\mu$s on the GPU.
NICP on template compares registered point sets ($4{,}096\times3$) and needs
811\,ms for the same search; ArcFace (512-dimensional cosine) needs
0.88\,ms on one CPU thread. Pairwise pipelines have no enrolment: a 1:$N$
search costs $N$ registrations, which at the measured per-pair cost amounts to
452\,s for similarity ICP with Chamfer and 3.6\,h for ICP with NICP at
$N=10{,}000$ (Table~\ref{tab:cost}).
% fonte: aau/runs/indomain_recog/summary.md, "ricerca su galleria iscritta (per query)" e "a coppie (per coppia)"

\begin{table}[!ht]
\centering
\small
\caption{Cost of enrolment and of 1:$N$ search per query (median). CPU
timings on one thread. Pairwise pipelines: per-pair time measured, 1:$N$
estimated as $N\times$ the per-pair median.}
\label{tab:cost}
% fonte: aau/runs/indomain_recog/summary.md, sez. "Tempi" (job 1060419, L40S + EPYC 9454) -- tranne la riga contrassegnata con $^\ddagger$
% fonte riga $^\ddagger$: aau/runs/evidence/e10/E10.md, "Per classe" (V100, 9.4k vertici, k 128, B=1: 97 ms)
\begin{tabular}{lcccc}
\toprule
Method & Enrolment / mesh & 1:100 & 1:1{,}000 & 1:10{,}000 \\
\midrule
\multicolumn{5}{l}{\emph{Enrolment methods}} \\
Ours, operators (CPU) + forward (GPU)   & 2.08\,s       & 19.8\,$\mu$s & 188\,$\mu$s & 1.96\,ms \\
Ours, operators (GPU)$^\ddagger$ + forward (GPU) & $\approx$0.15\,s \todo{end-to-end} & 40.2\,$\mu$s$^\ast$ & 50.0\,$\mu$s$^\ast$ & 68.0\,$\mu$s$^\ast$ \\
NICP on template (CPU)                  & 1.15\,s       & 7.33\,ms     & 81.2\,ms    & 811\,ms \\
ArcFace, 3 rendered views (CPU)         & 0.57\,s       & 8.1\,$\mu$s  & 76.3\,$\mu$s & 0.88\,ms \\
\midrule
\multicolumn{5}{l}{\emph{Pairwise methods} (per-pair cost in parentheses)} \\
Similarity ICP + Chamfer (45.2\,ms)     & --            & 4.5\,s       & 45\,s       & 452\,s \\
Similarity ICP + NICP P2Tri (1.29\,s)   & --            & 129\,s       & 1{,}286\,s  & 3.6\,h \\
\bottomrule
\end{tabular}

\vspace{0.3em}
\footnotesize{$^\ddagger$ Operators measured on a V100 for a 9.4k-vertex mesh; forward pass (55\,ms) on the L40S sample. Not an end-to-end measurement. $^\ast$ Search on the GPU.}
\end{table}

\paragraph{Reading.}
For one-off comparisons of a few meshes, cost does not decide anything: a
pairwise rigid ICP with Chamfer costs 34--45\,ms per pair.
% fonte: aau/runs/indomain_recog/summary.md, "a coppie (per coppia)" (ICP rigido + Chamfer 33.60 ms; ICP di similarita' + Chamfer 45.24 ms)
The difference appears in evaluation and retrieval at scale. At
$N=10{,}000$, counting the enrolment of the query, one search costs about
2\,s with NICP on template (1.15\,s + 0.81\,s), 0.57\,s with ArcFace, and
about 0.15\,s with our encoder on the GPU, against 452\,s for pairwise
similarity ICP with Chamfer and 3.6\,h for pairwise NICP.
% fonte: somme delle righe di Tab. tab:cost (aau/runs/indomain_recog/summary.md; operatori GPU da aau/runs/evidence/e10/E10.md); 0.15 s non e' end-to-end, vedi \todo sopra
Among enrolment methods, ours and ArcFace search fastest, but
cost must be read together with accuracy: on HIFI3D, NICP on template
reaches a rank-1 of 0.875 against 0.782 for our earlier encoder, and pairwise
rigid ICP with Chamfer reaches 0.996.
% fonte: aau/runs/competitors_hifi3d/summary.md, "Riconoscimento d'identita'", PRIMARIO (NICP su template 0.875; Rigid ICP + Chamfer 0.996); aau/runs/data_scale_ood/arcface_vs_scale_hifi3d.md, (b) (e108 0.782)
The accuracy of the final model at this cost is \todo{rank-1 e Spearman FR del modello finale su HIFI3D, FaMoS test, NoW}.
